\section{Flattening genuine nested section opens}
\begin{lemma}[Actual base and product coordinate identities]
\label{hir:actual-nested-section-coordinate-equivalences}
\lean{TranscendentalNakamaye.Pending.HironakaActualNestedSectionCoordinates.globalImage,TranscendentalNakamaye.Pending.HironakaActualNestedSectionCoordinates.baseEquiv,TranscendentalNakamaye.Pending.HironakaActualNestedSectionCoordinates.productWindow,TranscendentalNakamaye.Pending.HironakaActualNestedSectionCoordinates.productEquiv,TranscendentalNakamaye.Pending.HironakaActualNestedSectionCoordinates.windowSection,TranscendentalNakamaye.Pending.HironakaActualNestedSectionCoordinates.windowProjection}
\leanok
\uses{hir:actual-mask-section-whole-representative}
Let U be an open subset of complex n-space and W an open subset of U.
Its image W0 under the open inclusion is open in complex n-space.
The natural map W to W0 is a biholomorphism, as is the natural map from
C=p_U^{-1}(W), regarded inside p^{-1}(U), to p^{-1}(W0).
\end{lemma}
\begin{proof}
The image W0 is open because an open subspace inclusion is open. Each map
forgets one subtype and leaves the underlying affine point unchanged. A point
of W0 has the unique corresponding point of U and belongs to W; this gives
the inverse base map. A point of p^{-1}(W0) has projection in W0, hence in U,
and its corresponding U-point belongs to W. This gives the inverse product
map. Both compositions are the identity on the underlying affine points,
hence on the subtypes. Each direction is holomorphic: after composing with
the actual open inclusion its coordinate functions are the coordinate identity.
The local holomorphicity criterion for open inclusions then proves the claim.
\end{proof}
\begin{lemma}[The same literal section and projection squares]
\label{hir:actual-nested-section-coordinate-squares}
\lean{TranscendentalNakamaye.Pending.HironakaActualNestedSectionCoordinates.section_square,TranscendentalNakamaye.Pending.HironakaActualNestedSectionCoordinates.projection_square}
\leanok
\uses{hir:actual-nested-section-coordinate-equivalences}
With these two genuine biholomorphisms e and b, write sC for the section W to C
and pC for its actual projection. Then b composed with sC is the section of W0
composed with e, and e composed with pC is its projection composed with b.
\end{lemma}
\begin{proof}
For the section square both underlying affine values are the same literal
linear zero inclusion applied to the original n-vector. For the projection
square both are the same literal linear projection applied to the original
(n+1)-vector. Equality of subtype elements follows from equality of these
values; no inverse chart outside its target is used.
\end{proof}
\begin{lemma}[The entire section carrier is retained]
\label{hir:actual-nested-section-coordinate-carrier}
\lean{TranscendentalNakamaye.Pending.HironakaActualNestedSectionCoordinates.carrier_iff}
\leanok
\uses{hir:actual-nested-section-coordinate-squares}
For any subset A of W, the set of z in C with pC(z) in A and last normal
coordinate zero corresponds under b exactly to the lifted section carrier
of e(A) in p^{-1}(W0).
\end{lemma}
\begin{proof}
The preceding projection square identifies the new projected point with e(pC(z)).
Because e is injective, membership of this point in e(A) is equivalent to
pC(z) belonging to A. The last normal coordinate is unchanged because b
leaves the entire affine vector unchanged. Conjoining these two equivalences
gives the whole carrier equality, with no replacement by a cylinder.
\end{proof}
